\noindent\textbf{\specname} \hfill {\small\specnote}\par\medskip
Ein Zustand ist ein positives, normiertes lineares Funktional; im spurklassigen Fall
schreiben wir ihn als Dichteoperator $\rho\geq 0$ mit $\operatorname{tr}\rho = 1$. Für einen
Effekt $0\leq E\leq \mathbb{1}$ liefert die Bornsche Regel die Wahrscheinlichkeit
$p = \operatorname{tr}(\rho E)$. Die Energie des Wasserstoffatoms ist „gequantelt“, $E_n=-13{,}6\,\mathrm{eV}/n^2$.
\begin{equation}
  \Bigl(-\frac{\hbar^2}{2\mu}\nabla^2-\frac{e^2}{4\pi\varepsilon_0 r}\Bigr)\psi_{n\ell m}
  = E_n\,\psi_{n\ell m},\qquad
  \int_{\mathbb{R}^3}\lvert\psi_{n\ell m}(\mathbf{x})\rvert^2\,\mathrm{d}^3x = 1 .
\end{equation}
\begin{equation}
  [L_j,L_k]=\mathrm{i}\hbar\,\varepsilon_{jkl}L_l,\quad
  \langle\varphi\vert\mathrm{e}^{-\mathrm{i}Ht/\hbar}\vert\psi\rangle,\quad
  Z=\sum_{k}\mathrm{e}^{-\beta E_k},\quad
  \mathcal{H}=\bigoplus_{N=0}^{\infty}\mathcal{H}^{\otimes_s N},\quad
  \mathfrak{su}(2)\cong\mathfrak{so}(3).
\end{equation}
Inline: $\Delta x\,\Delta p\geq\hbar/2$, $\sigma_x\sigma_y=\mathrm{i}\sigma_z$, $\partial_t\rho=-\tfrac{\mathrm{i}}{\hbar}[H,\rho]$, $\alpha\beta\gamma\delta\epsilon\vartheta\kappa\lambda\xi\phi\varphi\omega\,\Gamma\Lambda\Omega$, $x^{2}_{i}$, $f'(x)$.
